
\newpage
\appendix
\onecolumn

\section{Appendix}\label{sec:appendix}

    %\subsection{Norm over edges} \label{sec:appendix_edgeNorm}
    %    For counting edges in \(\mathcal{KG}\) \emph{as if} they were undirected, we define
    %    \[
    %        ||_U = \sum_{(h,\dots,t)\in \mathcal{E}}
    %        \begin{cases} 
    %            0.5, & \text{if } (h,\dots,t),(t,\dots,h)\in \mathcal{E}, \\
    %            1,   & \text{otherwise}.
    %        \end{cases}
    %    \]

    \subsection{Sketch of Proof of Lemma 1} \label{sec:proof_lemma1}
        \begin{proof}[Proof]
            We break the argument into three parts.
            \paragraph{1. Counting all potential $(n+1)$-node sequences.}

                A simple \emph{n-hop path} is determined by choosing an ordered tuple of $(n+1)$ distinct entities. The number of ways to choose \emph{distinct} nodes $v_0,v_1,...,v_n$ out of $|\mathcal{V}|$ is:
                \[
                \binom{|\mathcal{V}|}{n+1} (n+1)!.
                \]
                
                Here, $\binom{|\mathcal{V}|}{n+1}$ is the number of ways to pick $n+1$ distinct nodes (unordered), and $(n+1)!$ is the number of ways to \emph{order} those nodes into a possible directed path of length $n$.

            \paragraph{2. Probability of each directed path being valid.}

                Under the random-graph assumption, the probability that there is a directed edge from $v_{i-1}$ to $v_i$ (for $i=1,...,n$) is $\frac{b}{|\mathcal{V}|-1}$.
                Assuming independence across edges, the probability that \emph{all} $n$ edges are present simultaneously is
                \[
                \Bigl(\tfrac{b}{|\mathcal{V}|-1}\Bigr)^{n}
                \]

            \paragraph{3. Expected number of valid $n$-hop paths.}

                By linearity of expectation (applied to each of the $\binom{|\mathcal{V}|}{n+1} (n+1)!$ ordered node tuples), the expected value of $|\mathcal{F}_n|$ is:

                \[ 
                \mathbb{E}[|\mathcal{F}_n|] = \binom{|\mathcal{V}|}{n+1} (n+1)! \Bigl(\tfrac{b}{|\mathcal{V}|-1}\Bigr)^{n}.
                \]

                For large $|\mathcal{V}|$, one can use the binomial approximation $\binom{n}{k}\leq\frac{n^k}{k!}$ for $\binom{|\mathcal{V}|}{n+1}$, such as

                \[
                \binom{|\mathcal{V}|}{n+1} \leq \frac{|\mathcal{V}|^{n+1}}{(n+1)!}
                \;\;\text{(for fixed $n$ as $|\mathcal{V}| \to \infty$ )}
                \]

            \paragraph{Computing Bound}

                By definition \ref{def:branching_factor} and \ref{def:phi_r},
                    \[
                    \phi_{n,r}\;=\;\frac{|\mathcal{F}_{n,r}|}{|\mathcal{F}_{A,r}|}
                    \;\;\text{and}\;\;
                    |\mathcal{F}_{A,r}| 
                    \;=\; |\mathcal{V}|\cdot b_r.
                    \]
                    
                Together with the result from the third step being $\mathbb{E}[|\mathcal{F}_n|] = \binom{|\mathcal{V}|}{n+1} (n+1)! \Bigl(\tfrac{b}{|\mathcal{V}|-1}\Bigr)^{n}$, we obtain,

                \begin{align*}
                    \mathbb{E}[\phi_{n,r}] &= \frac{|\mathcal{F}_{n,r}|}{|\mathcal{V}|\cdot b_r} \\
                    &= \frac{\binom{|\mathcal{V}|}{n+1} (n+1)! \Bigl(\tfrac{b_r}{|\mathcal{V}|-1}\Bigr)^{n}}{|\mathcal{V}|\cdot b_r}\\
                    &\leq \frac{\frac{|\mathcal{V}|^{n+1}}{(n+1)!} (n+1)! b_r^n}{|\mathcal{V}|b_r(|\mathcal{V}|-1)^n} \\
                    &= \frac{ |\mathcal{V}|^n b_r^{n-1}}{(|\mathcal{V}|-1)^n} \\
                    &= b_r^{n-1} \Bigl(\frac{ |\mathcal{V}| }{(|\mathcal{V}|-1)}\Bigr)^n\\
                    &= b_r^{n-1} \Bigl(\frac{ 1 }{1-\frac{1}{|\mathcal{V}|}}\Bigr)^n\\
                \end{align*}

                Which gives us the asymtotic upper bound,

                    \[
                    \lim_{|\mathcal{V}|\to\infty} \mathbb{E}[\phi_{n,r}] = \lim_{|\mathcal{V}|\to\infty} b_r^{n-1} \Bigl(\frac{ 1 }{1-\frac{1}{|\mathcal{V}|}}\Bigr)^n \to b_r^{n-1} 
                    \]

                Hence $\phi_{n,r}$ (and therfore $\phi_n$ overall) remains \emph{boudned above} by $b^{n-1}$, completing the proof.

        \end{proof}



    \subsection{Proof of Lemma 2} \label{sec:proof_lemma2}
        \begin{proof}[Sketch of Proof]
            By definition \ref{def:branching_factor} and \ref{def:phi_r},
                    \[
                    \phi_{n,r}\;=\;\frac{|\mathcal{F}_{n,r}|}{|\mathcal{F}_{A,r}|}
                    \;\;\text{and}\;\;
                    |\mathcal{F}_{A,r}| 
                    \;=\; |\mathcal{V}|\cdot b_r.
                    \]
                    
            From \ref{sec:proof_lemma1}, we have
            \[
            |\mathcal{F}_{n,r}| 
            \;\lesssim\;
            \binom{|\mathcal{V}|}{n+1}(n+1)!\,\Bigl(\frac{b_r}{|\mathcal{V}|-1}\Bigr)^n.
            \]
            
            because $\frac{b_r}{|\mathcal{V}|-1}$ is the approximate probability that a randomly chosen edge belongs to relation $r$.

            Thus,
            \begin{align*}
                \phi_{n,r} &= \frac{|\mathcal{F}_{n,r}|}{|\mathcal{F}_{A,r}|} =\frac{|\mathcal{F}_{n,r}|}{|\mathcal{V}|\,b_r} \\
                & \lesssim \binom{|\mathcal{V}|}{n+1}\frac{(n+1)!\,b_r^{\,n-1}}{|\mathcal{V}|\,(|\mathcal{V}|-1)^n}.
            \end{align*}
            
            \textbf{Full generalization} requires $\phi_{n,r}\ge\phi_G$ for \emph{all} $r$ yielding the required constraint
            \begin{align*}
            \phi_G &\leq\phi_{n,r} \\
            \Leftrightarrow \phi_G&\leq \binom{|\mathcal{V}|}{n+1}\frac{(n+1)!\,b_r^{\,n-1}}{|\mathcal{V}|\,(|\mathcal{V}|-1)^n} \\
            \Leftrightarrow b_r &\geq \sqrt[n-1]{
            \frac{\phi_G\,|\mathcal{V}|\,(|\mathcal{V}|-1)^n}{\binom{|\mathcal{V}|}{n+1}\,(n+1)!}
            }.
            \end{align*}
            
            In other words, if \emph{any} relation $r$ has $b_r$ (its branching factor) that fails to exceed this threshold, then $\phi_{n,r}$ cannot reach $\phi_G$. Hence, that particular relation will not “grok,” so the $KB$ is \emph{not fully generalizable} over $n$-hop facts (although it might still be \emph{partially} generalizable for other relations).
        \end{proof}


    \subsection{Formal Derivation of the Node-Count Bound}\label{sec:appendix_boundNlower}

        Similar to \ref{sec:proof_lemma2}, for \textbf{full generalization} we can derive a bound for the node count $|\mathcal{V}|$. From the requirement $\forall r \in \mathcal{R}, \phi_{n,r}\ge\phi_G$, we derive,
            \begin{align*}
                \forall r \in \mathcal{R}, \phi_{n,r}&\ge\phi_G \\
            \Leftrightarrow \min_r \binom{|\mathcal{V}|}{n+1}\frac{ (n+1)!b_r^{n-1}}{|\mathcal{V}|(|\mathcal{V}|-1)^n} &\geq \phi_G
            \end{align*}
                
        in other words, the worst generalizable relation $r$ still needs to fulfill $\phi_{n,r} \geq \phi_G $ for $KB$ to be fully generalizable.
        Resolving after $|\mathcal{V}|$ without the use of the binomial approximation yields,
        
        \begin{align*}
            \min_r \binom{|\mathcal{V}|}{n+1}\frac{ (n+1)!b_r^{n-1}}{|\mathcal{V}|(|\mathcal{V}|-1)^n} &\geq \phi_G \\
            \Leftrightarrow \min_r \frac{|\mathcal{V}|!}{(n+1)!(|\mathcal{V}|-n-1)!}\frac{ (n+1)!b_r^{n-1}}{|\mathcal{V}|(|\mathcal{V}|-1)^n} &\geq \phi_G \\
            \Leftrightarrow \min_r \frac{|\mathcal{V}|!}{(|\mathcal{V}|-n-1)!}\frac{ b_r^{n-1}}{|\mathcal{V}|(|\mathcal{V}|-1)^n} &\geq \phi_G \\
            \Leftrightarrow \frac{(|\mathcal{V}|-1)!}{(|\mathcal{V}|-n-1)!(|\mathcal{V}|-1)^n} \min_r b_r^{n-1} &\geq \phi_G \\
            \Leftrightarrow \frac{(|\mathcal{V}|-1)!}{(|\mathcal{V}|-n-1)!(|\mathcal{V}|-1)^n}   &\geq \max_r \frac{\phi_G}{b_r^{n-1}} \\
            \Leftrightarrow \frac{\Gamma(|\mathcal{V}|)}{\Gamma(|\mathcal{V}|-n)(|\mathcal{V}|-1)^n}   &\geq \max_r \frac{\phi_G}{b_r^{n-1}}
        \end{align*}


            Thus we have:
            {\small
            \[
            |\mathcal{V}|
            \;\ge\;
            \min\Bigl\{
               v \in \mathbb{N}
               : 
               \frac{\Gamma(v)}{\Gamma(v-n\bigr)\,(v-1)^n}
               \;\ge\;\max_r \frac{\phi_G}{b_r^{n-1}}
            \Bigr\}.
            \]
            }

        %Above, we have derived an exact bound depending on the $\Gamma$ function.
        For an empirical example (i.e., using a randomly generated graph), see Figure \ref{fig:A_phi_n3_b2}.

        \begin{figure}[ht]
            \centering
            \includesvg[width=0.45\textwidth]{phi_n3_b2.svg}
            \caption{
                Growth of $\phi_{3,r}$ (\textbf{y-axis}) with $|\mathcal{V}|$ (\textbf{x-axis}) for $b_r=2$.
                The \textbf{red line} is $\phi_{3,r}=\binom{|\mathcal{V}|}{4}\frac{96}{|\mathcal{V}|(|\mathcal{V}|-1)^3}$ (formula values). 
                The \textbf{green line} are empirical values obtained by randomly generating a graph with the same amount of nodes ($|\mathcal{V}|$) and edges ($|\mathcal{V}| b_r$).
                Due to randomness, our generated/empirical graph can have locally higher branching factors, resulting in overall more inferred facts. Our formula for $\phi_{3,r}$ therefore effectively underestimates the true value of $\phi_{3,r}$. Nevertheless, the formula correctly approximates the shape and order of magnitude. This holds also for other combinations of $b_r$ and $n$.
                }
            \label{fig:A_phi_n3_b2}
        \end{figure}


    \subsection{Qualitative Examples}\label{sec:appendix_qualitative_examples}

    In the following section, we present qualitative examples of QA pairs (\textit{"Question", "Ground Truth"}) from the composition and comparison tasks that the model failed to classify correctly. Some errors stem from inconsistencies in the 2WikiMultiHopQA dataset, while others arise due to our augmentation strategy. Overall, we believe that the model is capable of achieving 100\% accuracy, as demonstrated in previous studies.
    
    \textbf{Composition}
    There are several grammatical inconsistencies in the 2WikiMultiHopQA dataset that prevent our model from reaching 100\% accuracy.
    
    \paragraph{Nationality questions may have inconsistent ground truth formats}
    For nationality-related questions, the ground truth can be either an adjective or the name of a country, even when the latter is grammatically incorrect.
    
    \begin{quote}
    \textbf{Question:} What nationality is William Seymour, 3rd Duke of Somerset's father? \\ 
    \textbf{Ground Truth:} English
    \end{quote}
    
    \begin{quote}
    \textbf{Question:} What nationality is Amadeus VIII, Duke of Savoy's mother? \\ 
    \textbf{Ground Truth:} \textbf{France}
    \end{quote}
    
    \paragraph{Adjective instead of noun in country-related answers}
    In some cases, the ground truth is an adjective instead of a noun referring to a country.
    
    \begin{quote}
    \textbf{Question:} Which country is the trumpeter of Paper Bird from? \\ 
    \textbf{Ground Truth:} American
    \end{quote}
    
    \textbf{Comparison}
	Comparison errors are primarily due to limitations in the data augmentation algorithm, which failed to generate a sufficient number of inferred facts for some relational queries. Most of them are related to lakes, rivers, or airports, which we attribute to the sparsity of the base data.

            
        \begin{quote}
        \textbf{Question:} Are both Wainwright/Wainwright Field 21 Airport and Roberval Air Saguenay Water Aerodrome located in the same country? \\ 
        \textbf{Ground Truth:} Yes
        \end{quote}
        
        \begin{quote}
        \textbf{Question:} Are Long Lake, East Ferris, Ontario, and Montreal Lake, Saskatchewan, both located in the same country? \\ 
        \textbf{Ground Truth:} Yes
        \end{quote}
        
        \begin{quote}
        \textbf{Question:} Are Deer Creek, Osage River, and Big Prairie Dog Creek both located in the same country? \\ 
        \textbf{Ground Truth:} Yes
        \end{quote}
        
        \begin{quote}
        \textbf{Question:} Are Chicoutimi/Saint-Honoré Aerodrome and Rtishchevo Air Base both located in the same country? \\ 
        \textbf{Ground Truth:} No
        \end{quote}

    
    \subsection{Data Synthesis}\label{sec:appendix_datasynthesis}

    Here, we present the prompts used throughout our data augmentation process. During our experiments, we leveraged both GPT-4o and o1-mini to generate diverse and high-quality augmented data.
    
        \subsubsection{Composition}
    
            \textbf{Graph parsing}
            
            \noindent
            \fbox{%
                \parbox{\textwidth}{%
                    \small\ttfamily
            You are graph gpt. You build graph based on the provided text. 
            Find all objects, their relations and types.
            
            Pick one of the following types:
            - Person
            - Location
            - Object (include everything that was not above)
            
            Return the following format with numbering: 
            1. <Avatar; Film><director><James Cameron; Person>
            2. <James Cameron; Person><directed><Titanic; Object>
            }
            }
        
            \textbf{Question formatting}
            
            \noindent
            \fbox{%
                \parbox{\textwidth}{%
                    \small\ttfamily
            You are a question formatting assistant. Your task is to create questions based on the given relations and objects. 
            
            Use the provided examples as a guide for the question style. Ensure that the answer remains unchanged and enclosed in <a> tags. 
            You may rephrase one question, given the example format. Strictly follow the logic of given examples. 
            Connect it in the following logic: <obj1> -> <rel1> -> <rel2> -> <obj3>
            
            Return numbered responses in format:
            1. What is the director of the film that James Cameron produced?<a>Steven Spielberg</a>
            2. Who directed the movie starring Tom Cruise?<a>Christopher Nolan</a>
            }
            }
        
        \subsubsection{Comparison}
    

            \textbf{Atomic fact generation}
            
            \noindent
            \fbox{%
                \parbox{\textwidth}{%
                    \small\ttfamily
            You are a helpful assistant that generates geographical facts.
                Generate new unique locations and their countries in the following format:
                Follow the style of the examples, but do not use the same locations.
                
                Rules:
                1. Use real locations and countries
                2. Each location should be unique
                3. DO NOT REUSE PROVIDED EXAMPLES
                4. Do not answer the question - only provide locations
                5. Do not use formatting except for numbering
                6. Generate equal amount of NEW!!! locations for following countries: {}
            
            }
            }
            
            \textbf{Detailed atomic fact generation}
            
            \noindent
            \fbox{%
                \parbox{\textwidth}{%
                    \small\ttfamily
            You are a helpful assistant that generates geographical facts.
                Based on the provided examples, generate a paragraph for each location-country pair. Strictly follow the style and lenght of the provided examples Do not answer the question - only provide the paragraph with numbering. DO not return empty lines. One by one. Return the number according to the given data. Here are the examples: 
                {}
            }
            }
            
            